\section{Video Experiments on KAIST}
\label{sec:video}
This section asks two questions about the video extension of Section~\ref{sec:tsm}. Does the fusion block help on a pedestrian benchmark with a long history of published results, and what does the Temporal State Memory add on top of a still-image detector, on clean video and while one camera fails? Every number here comes from one training run per model and one memory training run, so no seed spread is available, and we say so wherever a difference is small.

\subsection{Protocol}
\label{sec:kaistprotocol}
KAIST~\cite{hwang2015kaist} supplies visible and thermal video at 20 frames per second. We use the standard split: sets 00 to 05 for training and sets 06 to 11 for testing, with sets 00 to 02 and 06 to 08 recorded by day and the rest at night. The still-image detectors train on every \cfgKaistTrainStride{}th frame of the training sets, with persons at least \cfgKaistTrainMinH{} pixels tall as targets. Both detectors, CFFM-Net v2 and the concat baseline with modality dropout, follow the recipe of Section~\ref{sec:experiments} with \cfgKaistEpochs{} epochs and seed zero. The test frames are every \cfgKaistTestStride{}th frame of the test sets, \kstNAll{} frames of which \kstNDay{} are by day and \kstNNight{} by night.

The metric is the log-average miss rate (MR) over nine false-positives-per-image values between $10^{-2}$ and $10^{0}$ on the reasonable subset: persons at least \cfgKaistMinH{} pixels tall and at most partly occluded (visible fraction of at least \cfgKaistMinVis{}). Lower is better. We use the original KAIST annotation files for training and testing and our own implementation of the Caltech matching protocol, not the KAIST evaluation toolbox, and we have not applied the sanitised annotations~\cite{li2018kaistsanitized} or the improved test annotations used by most recent papers. The miss rates below are therefore internally consistent but are not directly comparable with published figures, and Section~\ref{sec:sota} places them next to those figures with this caveat.

For video the detector runs causally at \cfgFailFps{} frames per second (every second frame) and is scored on the same test grid. A failure schedule repeats \cfgFailSeconds{} healthy seconds and \cfgFailSeconds{} failed seconds. We report the miss rate on all frames and on the failed frames alone. Conditions are a lost thermal camera (a black image), a lost thermal camera declared by the sensor flag, a darkened visible camera, and additive noise on the thermal camera, all at severity three. Detector weights are identical in every row; only the memory is switched on or off at test time. Tracking uses ByteTrack at the full 20 frames per second with a detection threshold of \cfgTrackConf{}, on four test sequences (two by day, two by night, \trkUniqueTotal{} annotated persons), scored by MOTA and IDF1.

\subsection{Still images}
\label{sec:kaiststills}
\begin{table}[t]
\centering
\caption{KAIST test results on single frames, reasonable subset. Lower is better for the miss rate. One training run per model, \kstNAll{} test frames, so no confidence interval is available.}
\label{tab:kaist}
{\footnotesize\setlength{\tabcolsep}{4pt}
\begin{tabular}{lccc}
\toprule
Measure & Concat & CFFM-Net v2 & Change \\
\midrule
Miss rate, all (\%) & \kstConcatAllMr & \kstCffmAllMr & \kstMarginAllMr \\
Miss rate, day (\%) & \kstConcatDayMr & \kstCffmDayMr & \kstMarginDayMr \\
Miss rate, night (\%) & \kstConcatNightMr & \kstCffmNightMr & \kstMarginNightMr \\
AP$_{50}$, all (\%) & \kstConcatAllApFifty & \kstCffmAllApFifty & \kstMarginAllApFifty \\
\bottomrule
\end{tabular}}
\end{table}

Table~\ref{tab:kaist} gives the single-frame results. CFFM-Net v2 does not beat the concat baseline on KAIST. Over all test frames the miss rate is \kstCffmAllMr{}\,\% against \kstConcatAllMr{}\,\%, a change of \kstMarginAllMr{}\,points that favours the baseline, and the AP$_{50}$ is \kstCffmAllApFifty{}\,\% against \kstConcatAllApFifty{}\,\%. By day the baseline is ahead by \kstMarginDayMr{}\,points of miss rate; at night the two are level (\kstMarginNightMr{}\,points). With one run each and no estimate of run-to-run spread, which on the detection benchmarks of Section~\ref{sec:main} was about one point of mAP, these differences are best read as a tie with a slight lean to the baseline. We make no claim that the reliability-gated scan improves still-image pedestrian detection on KAIST. Absolute miss rates are high for the reasons in Section~\ref{sec:kaistprotocol} and because both models are nano scale and trained for \cfgKaistEpochs{} epochs on a subsample.

\subsection{Temporal State Memory}
\label{sec:kaistmemory}
The memory is trained for \cfgMemEpochs{} epochs on clips of \cfgClipLen{} frames with the detector frozen (batch \cfgMemBatch{}, learning rate \cfgMemLr{}), and is attached to the CFFM-Net v2 trained on KAIST. All rows of Table~\ref{tab:memory} use that one memory.

\begin{table*}[t]
\centering
\caption{Temporal State Memory against the same detector without memory on KAIST. A camera fails for four seconds in every eight. Failed frames are those inside a failure. Change is memory minus image. One memory training run, one detector.}
\label{tab:memory}
{\footnotesize\setlength{\tabcolsep}{4pt}
\begin{tabular}{lcccccc}
\toprule
& \multicolumn{3}{c}{All frames} & \multicolumn{3}{c}{Failed frames} \\ \cmidrule(lr){2-4}\cmidrule(lr){5-7} Condition & Image & Memory & Change & Image & Memory & Change \\
\midrule
\multicolumn{7}{l}{\emph{Miss rate (\%), lower is better}} \\
Clean & \memCleanImageAllMr & \memCleanMemoryAllMr & \memChangeCleanAllMr & -- & -- & -- \\
Thermal lost & \memThermalLostImageAllMr & \memThermalLostMemoryAllMr & \memChangeThermalLostAllMr & \memThermalLostImageFailedMr & \memThermalLostMemoryFailedMr & \memChangeThermalLostFailedMr \\
Thermal lost, flagged & \memThermalLostFlaggedImageAllMr & \memThermalLostFlaggedMemoryAllMr & \memChangeThermalLostFlaggedAllMr & \memThermalLostFlaggedImageFailedMr & \memThermalLostFlaggedMemoryFailedMr & \memChangeThermalLostFlaggedFailedMr \\
Visible dark & \memVisibleDarkImageAllMr & \memVisibleDarkMemoryAllMr & \memChangeVisibleDarkAllMr & \memVisibleDarkImageFailedMr & \memVisibleDarkMemoryFailedMr & \memChangeVisibleDarkFailedMr \\
Thermal noise & \memThermalNoiseImageAllMr & \memThermalNoiseMemoryAllMr & \memChangeThermalNoiseAllMr & \memThermalNoiseImageFailedMr & \memThermalNoiseMemoryFailedMr & \memChangeThermalNoiseFailedMr \\
\midrule
\multicolumn{7}{l}{\emph{Tracking, mean of four sequences}} \\ & \multicolumn{3}{c}{MOTA (\%)} & \multicolumn{3}{c}{IDF1 (\%)} \\ \cmidrule(lr){2-4}\cmidrule(lr){5-7} Condition & Image & Memory & Change & Image & Memory & Change \\
Clean & \trkCleanImageMota & \trkCleanMemoryMota & \trkChangeCleanMota & \trkCleanImageIdfOne & \trkCleanMemoryIdfOne & \trkChangeCleanIdfOne \\
Thermal lost & \trkThermalLostImageMota & \trkThermalLostMemoryMota & \trkChangeThermalLostMota & \trkThermalLostImageIdfOne & \trkThermalLostMemoryIdfOne & \trkChangeThermalLostIdfOne \\
\bottomrule
\end{tabular}}
\end{table*}


On clean video the memory changes the miss rate over all frames from \memCleanImageAllMr{}\,\% to \memCleanMemoryAllMr{}\,\% (\memChangeCleanAllMr{}\,points) and AP$_{50}$ by \memChangeCleanAllApFifty{}\,points. This is a gain of half a point from one run, and we do not treat it as established.

Under failures the picture depends on the fault. When the thermal camera is lost the memory does not help where it should: over all frames the miss rate moves by \memChangeThermalLostAllMr{}\,points, but on the failed frames it is \memThermalLostMemoryFailedMr{}\,\% with memory against \memThermalLostImageFailedMr{}\,\% without (\memChangeThermalLostFailedMr{}\,points, slightly worse), and at night the all-frame result is also worse (\memChangeThermalLostNightMr{}\,points). Declaring the lost camera through the sensor flag leaves this unchanged (\memChangeThermalLostFlaggedFailedMr{}\,points on failed frames). The all-frame improvement comes from the frames outside the failures, where the miss rate falls from \memThermalLostImageHealthyMr{}\,\% to \memThermalLostMemoryHealthyMr{}\,\%. A lost thermal camera is the case a recurrent state should ride out, and a four-second outage appears to be too long for this memory to bridge, though we did not test other outage lengths. When the visible camera is darkened the memory helps on the failed frames, \memVisibleDarkImageFailedMr{}\,\% to \memVisibleDarkMemoryFailedMr{}\,\% (\memChangeVisibleDarkFailedMr{}\,points), with the largest part of the gain at night (\memChangeVisibleDarkNightMr{}\,points over all night frames and \memChangeVisibleDarkDayMr{} by day).

The large effect is under thermal noise. Without memory, noise at severity three on the thermal camera drives the miss rate from \memCleanImageAllMr{}\,\% to \memThermalNoiseImageAllMr{}\,\%, and to \memThermalNoiseImageFailedMr{}\,\% on the failed frames; with memory it is \memThermalNoiseMemoryAllMr{}\,\% overall and \memThermalNoiseMemoryFailedMr{}\,\% on failed frames, a change of \memChangeThermalNoiseAllMr{} and \memChangeThermalNoiseFailedMr{}\,points, and AP$_{50}$ rises from \memThermalNoiseImageAllApFifty{}\,\% to \memThermalNoiseMemoryAllApFifty{}\,\%. This comparison has a confound that we cannot remove with the runs we have. The memory was trained on clips with injected faults, noise among them, while the detector it sits on was not, so the gain may come from the memory block learning to suppress noisy features rather than from carrying state across frames. A control that trains a stateless adapter of the same size on the same clips would separate the two, and we have not run it. The result shows that the memory block can recover most of the accuracy lost to thermal noise; it does not show that the recurrent state is the reason.

\subsection{Tracking}
\label{sec:kaisttracking}
The memory does not improve tracking in our runs (Table~\ref{tab:memory}). On clean video MOTA falls from \trkCleanImageMota{}\,\% to \trkCleanMemoryMota{}\,\% and IDF1 from \trkCleanImageIdfOne{}\,\% to \trkCleanMemoryIdfOne{}\,\%, and the memory is ahead on MOTA in \trkCleanMemoryWinsMota{} of \trkSeqCount{} sequences. With the thermal camera failing, the cost of the failure itself is large without memory (MOTA down by \trkFaultCostImageMota{}\,points), and memory recovers none of it: MOTA is \trkThermalLostMemoryMota{}\,\% against \trkThermalLostImageMota{}\,\%, IDF1 is \trkThermalLostMemoryIdfOne{}\,\% against \trkThermalLostImageIdfOne{}\,\%, and the memory wins MOTA on \trkThermalLostMemoryWinsMota{} of \trkSeqCount{} sequences. Identity switches are close (\trkCleanMemorySwitches{} against \trkCleanImageSwitches{} on clean video and \trkThermalLostMemorySwitches{} against \trkThermalLostImageSwitches{} with the failure, summed over the four sequences). Four sequences give a coarse estimate, but the sign of the IDF1 change is the same in both conditions. We therefore report the memory as a detection-time component whose benefit shows under visible darkness and thermal noise, and do not claim a tracking benefit.

\paragraph{Supplementary material}
Supplementary video S1 (\texttt{failure\_memory.mp4}, ten seconds at ten frames per second) shows a night sequence at the moment the thermal camera is cut, with the image-only detector above and the memory detector below.
